% Part: incompleteness
% Chapter: theories-computability
% Section: inseparability

\documentclass[../../../include/open-logic-section]{subfiles}

\begin{document}

\olfileid{inc}{tcp}{ins}

\olsection{$\Th{Q}$-তে প্রমাণযোগ্য ও খণ্ডনযোগ্য বাক্যগুলি গণনসাধ্যভাবে
  অবিচ্ছেদ্য}


$\Th{\bar Q}$ দিয়ে সেই সব বাক্যের সেট বোঝাই, যাদের \emph{অস্বীকার}
$\Th{Q}$-তে প্রমাণযোগ্য; অর্থাৎ,
$\Th{\bar Q} = \Setabs{!A}{\Th{Q} \Proves \lnot !A}$। মনে করো,
বিচ্ছিন্ন সেট~$A$ ও $B$-কে \emph{গণনসাধ্যভাবে অবিচ্ছেদ্য} বলা হয় যদি
এমন কোনো গণনসাধ্য সেট~$C$ না থাকে যাতে $A \subseteq C$ এবং
$B \subseteq \Complement{C}$।

\begin{lem}
$\Th{Q}$ ও $\Th{\bar Q}$ গণনসাধ্যভাবে অবিচ্ছেদ্য।
\end{lem}

\begin{proof}
ধরি $C$ এমন একটি গণনসাধ্য সেট যাতে $\Th{Q} \subseteq C$ এবং
$\Th{\bar Q} \subseteq \Complement{C}$। $R(x,y)$ নিচের সম্বন্ধটি হোক:
\begin{quote}
$x$ কোনো সূত্র~$!D(u)$-কে সংকেতায়িত করে এবং $!D(\num y)$ $C$-তে আছে।
\end{quote}
দেখাব যে $R(x,y)$ একটি সার্বজনীন গণনসাধ্য সম্বন্ধ; তাতেই বিরোধ পাওয়া
যাবে।

ধরি $S(y)$ গণনসাধ্য এবং $\Th{Q}$-তে $!D_S(u)$ দ্বারা উপস্থাপিত। তাহলে
\begin{eqnarray*}
S(n) & \lif & \Th{Q} \Proves !D_S(\num n) \\
& \lif & !D_S(\num n) \in C
\end{eqnarray*}
এবং
\begin{eqnarray*}
\lnot S(n) & \lif & \Th{Q} \Proves \lnot !D_S(\num n) \\
& \lif & !D_S(\num n) \in \Th{\bar Q} \\
& \lif & !D_S(\num n) \not\in C
\end{eqnarray*}
সুতরাং $S(y)$ সম্বন্ধটি $R(\Gn{!D_S(u)},y)$-এর সমতুল্য।
% BN-SRC-241 TARGET-CORRECTION-BEGIN
\par\smallskip\noindent\textnormal{\small\textbf{উৎস-সংশোধন (BN-SRC-241):}
হিমায়িত সমীকরণে মেটা-সম্বন্ধ~$S$-এর যুক্তি হিসেবে বস্তু-ভাষার সংখ্যাপদ
$\num n$ লেখা হয়েছে; আর শেষ বাক্যে সূত্রের কোড নিতে
$R(\#(!D_S(\num u)),y)$ লেখা হয়েছে, যেখানে মুক্ত চল~$u$-কে সংখ্যাপদে
রূপ দেওয়া অনুপযুক্ত। বাংলা পাঠে মেটা-স্তরে $S(n)$ এবং উপস্থাপক
সূত্রটির গ্যোডেল কোড $\Gn{!D_S(u)}$ রাখা হয়েছে।}
% BN-SRC-241 TARGET-CORRECTION-END
\end{proof}

\end{document}
